\chapter{Conclusion and Recommendation}
\label{ch:conclusion}

\section{Summary of Findings}

This project set out to determine whether \gls{gcn} circular text alone can support a fast, approximate chirp-mass estimate. The answer, established with evidence rather than assumed: not reliably with the standard \gls{snr}-scaling physics formula, because network \gls{snr} and orbital inclination are structurally absent from circular text -- confirmed both by a data-availability audit and by the formula's failure even after properly calibrating its one free constant (Chapter~\ref{ch:results}). A statistical alternative based on distance alone performs meaningfully better, though a trivial classification-based baseline tempers how much of that improvement is genuinely new information.

\section{Contribution}

The core contribution is not a proof that circular-level chirp-mass estimation works well -- it does not, cleanly, and this report states that plainly. The contribution is: a reproducible, provenance-tracked pipeline converting unstructured circular text into a structured, auditable dataset; an honest, quantified feasibility result that explains \emph{why} the physics-motivated approach falls short and what a data-driven alternative achieves instead; and a literature-grounded case, backed by a direct comparison study, for why an AI-assisted tool in this domain should be built on retrieval rather than model memory.

\section{Limitations}

\begin{itemize}
    \item The validation set is small (49 events); leave-one-out cross-validation reduces, but does not eliminate, overfitting risk at this sample size.
    \item Reference-mass events are, by construction, drawn from later observing runs than O4a, since O4a circulars never report this field; O4a data contributed to the availability audit but no validation examples.
    \item The statistical model uses distance alone; \gls{far} and classification probability were tested as additional predictors and did not help at this sample size, though this may change as the archive accumulates more circulars with reference chirp mass.
    \item No external reference catalog (e.g. \gls{gwtc}/\gls{gwosc}) was incorporated, by scope decision; doing so was identified as a possible future extension pending supervisor guidance on scope.
\end{itemize}

\section{Recommendation}

Given the coarse, bin-quantized nature of the reference chirp mass values circulars do report, this project recommends that future work on this problem adopt ``predict the correct mass bin'' rather than a fixed percentage-error target as the primary success criterion, since the latter implicitly demands more precision than the ground truth itself provides.

\section{Future Work}

\begin{itemize}
    \item Expand the historical validation sample as more circulars with reference chirp mass accumulate, or by incorporating an external catalog if scope permits.
    \item Test whether source classification, more carefully combined with distance (rather than as a simple alternative baseline), improves on the current statistical estimator.
    \item Develop an uncertainty-aware estimator rather than a point estimate.
    \item Complete the AI-model comparison study on a larger, fresh event batch, and complete a working (not just designed) GW Pathfinder prototype.
    \item Explore additional circular-derivable parameters beyond chirp mass (e.g. sky localization, detector count, event time) not yet incorporated into the pipeline.
    \item Prepare a publication-quality version of the statistical analysis for the strongest result, per the project's longer-term research-paper goal.
\end{itemize}
